% Wersja źródłowa: audyt II, jednostka 04
\setcounter{section}{3}
\section{Zbiory ograniczone, rodziny operatorów i metryzowalność TVS}\label{sec:rodziny-metryzowalnosc}

\subsection{Motywacja i przyjęte założenia}

W przestrzeni unormowanej liczba $\sup_{x\in B}\|x\|$ mierzy
rozmiar zbioru $B$. W ogólnej przestrzeni liniowo-topologicznej
nie mamy wyróżnionej normy. Możemy jednak określić, czy wszystkie
punkty zbioru $B$ mieszczą się w odpowiednim wielokrotnym
powiększeniu dowolnego otoczenia zera. W \hyperref[sec:tvs]{rozdziale 1} przyjęliśmy
następującą definicję: $B\subseteq X$ jest \emph{ograniczony}, gdy
dla każdego otoczenia $U$ zera istnieje $t>0$ takie, że
$B\subseteq tU$. Jest to ograniczoność w sensie von Neumanna.

W całym rozdziale $\K\in\{\R,\C\}$, a $X,Y$ oznaczają
przestrzenie liniowo-topologiczne nad $\K$; jeśli nie wskazano
inaczej, są one Hausdorffa. Otoczenia zera nie muszą być otwarte.
Przez $\mathcal L(X,Y)$ rozumiemy \emph{zbiór ciągłych odwzorowań
liniowych} $X\to Y$. Dla rodziny odwzorowań $\mathcal T$
zapis $\mathcal T(B)$ oznacza zbiór
$\{T(x):T\in\mathcal T,\ x\in B\}$.

\subsection{Sekwencyjne kryterium ograniczoności}

\begin{twierdzenie}[Kryterium ograniczoności przez skalarne ciągi]
\label{thm:bounded-sequential-test}
Niech $X$ będzie przestrzenią liniowo-topologiczną, a
$B\subseteq X$. Następujące warunki są równoważne:
\begin{enumerate}[label=\textup{(\roman*)}]
\item $B$ jest ograniczony w sensie von Neumanna;
\item dla każdego ciągu $(x_n)$ w $B$ i każdego ciągu
      $(\lambda_n)$ w $\K$ zbieżnego do zera zachodzi
      $\lambda_nx_n\to0$ w $X$.
\end{enumerate}
\end{twierdzenie}
\begin{proof}
Załóżmy (i). Niech $U$ będzie otoczeniem zera; wybierzmy
zrównoważone otoczenie $V\subseteq U$.
Istnieje $t>0$ takie, że $B\subseteq tV$.
Gdy $|\lambda_n|t\le1$, dostajemy
$\lambda_nx_n\in\lambda_ntV\subseteq V\subseteq U$.
Zatem $\lambda_nx_n\to0$.

Jeżeli (i) nie zachodzi, istnieje zrównoważone otoczenie
$U$ zera takie, że $B\not\subseteq tU$ dla żadnego $t>0$.
Dla każdego $n\ge1$ wybierzmy $x_n\in B\setminus nU$
i połóżmy $\lambda_n=1/n$. Wtedy $\lambda_n\to0$,
ale $\lambda_nx_n\notin U$ dla żadnego $n$, co przeczy (ii).
\end{proof}

\begin{stwierdzenie}[Podstawowe działania na zbiorach ograniczonych]
\label{prop:bounded-operations}
Niech $X$ będzie TVS, a $B,C\subseteq X$ zbiorami ograniczonymi.
Wówczas $B\cup C$, $B+C$, $\lambda B$ dla $\lambda\in\K$
oraz domknięcie $\overline B$ są ograniczone. Każdy podzbiór
zbioru ograniczonego jest ograniczony.
\end{stwierdzenie}
\begin{proof}
Wybierzmy zrównoważone otoczenie $V$ zera z
$V+V\subseteq U$, gdzie $U$ jest zadanym otoczeniem zera.
Istnieją liczby $s,t>0$ z $B\subseteq sV$ i $C\subseteq tV$.
Jeżeli $r=\max(s,t)$, to $B\cup C\subseteq rV$, ponieważ
$V$ jest zrównoważony. Ponadto
$B+C\subseteq sV+tV\subseteq (s+t)(V+V)\subseteq(s+t)U$;
wystarcza też użyć $r$ i sumy $V+V$. Ograniczoność
$\lambda B$ i podzbiorów wynika bezpośrednio z definicji.

Dla domknięcia wybierzmy otwarte zrównoważone otoczenie
$V$ z $V+V\subseteq U$. Zachodzi $\overline V\subseteq V+V$:
gdy $x\in\overline V$, otwarty zbiór $x+V$ przecina $V$,
zatem $x\in V-V=V+V$. Skoro $B\subseteq tV$ dla pewnego
$t>0$, mamy
$\overline B\subseteq t\overline V\subseteq tU$.
\end{proof}

\begin{uwaga}[Ważna pułapka: otoczka wypukła]
Dla $B\subseteq X$ przez $\operatorname{conv}(B)$ rozumiemy
zbiór wszystkich skończonych kombinacji wypukłych punktów $B$
(przyjmując $\operatorname{conv}(\varnothing)=\varnothing$).
W \emph{dowolnej} TVS otoczka wypukła zbioru ograniczonego
nie musi być ograniczona. Własność ta stanie się prawdziwa
po założeniu lokalnej wypukłości: dla dowolnego otoczenia
$U$ zera wybieramy wypukłe, zrównoważone otoczenie
$V\subseteq U$. Jeżeli $B\subseteq tV$, to
$\operatorname{conv}(B)\subseteq tV\subseteq tU$, więc
otoczka wypukła jest ograniczona.
\end{uwaga}

\begin{przyklad}[Nieograniczona otoczka wypukła kuli jednostkowej]
\label{ex:quasinorm-convex-hull}
Niech $0<p<1$ oraz $X=L^p([0,1])$ z metryką
$d(f,g)=\int_{[0,1]}|f(t)-g(t)|^p\,\mathrm d\lambda(t)$.
Zbiór $B=\{f:\int_{[0,1]}|f(t)|^p\,\mathrm d\lambda(t)\le1\}$ jest
ograniczony: dla dowolnego $\varepsilon>0$ jego odpowiednio
mała homotetia mieści się w
$\{f:\int_{[0,1]}|f(t)|^p\,\mathrm d\lambda(t)<\varepsilon\}$.
Dla $n\ge1$ podzielmy $[0,1]$ na przedziały $I_{n,j}$,
$j=1,\ldots,n$, o długości $1/n$ i połóżmy
$f_{n,j}=n^{1/p}\mathbf1_{I_{n,j}}$. Wtedy
$\int_{[0,1]}|f_{n,j}(t)|^p\,\mathrm d\lambda(t)=1$,
natomiast ich średnia arytmetyczna
$g_n=n^{-1}\sum_{j=1}^n f_{n,j}=n^{1/p-1}\mathbf1_{[0,1]}$
należy do $\operatorname{conv}(B)$ i spełnia
\[
 \int_{[0,1]}|g_n(t)|^p\,\mathrm d\lambda(t)=n^{1-p}\longrightarrow\infty.
\]
Otoczka wypukła $B$ nie jest więc ograniczona. Dowodzi to
również, że $X$ nie ma bazy otoczeń zera złożonej ze
zbiorów wypukłych: w przeciwnym razie otoczka wypukła
każdego zbioru ograniczonego byłaby ograniczona.
\end{przyklad}

\subsection{Lokalna ograniczoność i jej granice}

\begin{definicja}[Przestrzeń lokalnie ograniczona]
Przestrzeń liniowo-topologiczną $X$ nazywamy
\emph{lokalnie ograniczoną}, jeżeli istnieje ograniczone
otoczenie zera w $X$.
\end{definicja}

\begin{stwierdzenie}[Kryterium lokalnej ograniczoności]
\label{prop:locally-bounded}
Dla TVS $X$ następujące warunki są równoważne:
\begin{enumerate}[label=\textup{(\roman*)}]
\item $X$ ma ograniczone otoczenie zera;
\item istnieje otoczenie zera $U$ takie, że dla każdego
otoczenia zera $V$ istnieje $t>0$ z $U\subseteq tV$;
\item istnieje baza otoczeń zera postaci $\{tU:t>0\}$
dla pewnego zrównoważonego otoczenia $U$.
\end{enumerate}
\end{stwierdzenie}
\begin{proof}
Równoważność (i) i (ii) jest treścią definicji.
Załóżmy (i) i wybierzmy zrównoważone otoczenie
$U$ zawarte w ograniczonym otoczeniu. Jest ono również
ograniczone. Każdy zbiór $tU$, $t>0$, jest otoczeniem
zera, a ograniczoność $U$ daje: dla dowolnego otoczenia
$V$ istnieje $s>0$ takie, że $U\subseteq sV$, czyli
$s^{-1}U\subseteq V$. Dowodzi to (iii).
Implikacja (iii)$\Rightarrow$(i) wynika z warunku, że
dla danego otoczenia $V$ pewna homotetia $tU$ jest w $V$.
\end{proof}

\begin{przyklad}[Dwa odmienne przypadki]
Przestrzeń $\K^{\N}$ z topologią produktową nie jest
lokalnie ograniczona. Każde jej otoczenie zera zawiera
zbiór, w którym tylko skończenie wiele współrzędnych jest
ograniczonych; poza tym zbiorem współrzędnych można wybrać
współrzędną $j$ całkowicie swobodną. Takie otoczenie nie
mieści się w żadnym wielokrotnym powiększeniu zbioru
$\{x:|x_j|<1\}$.

Z kolei $L^p([0,1])$, $0<p<1$, jest lokalnie ograniczona:
metryczna kula zera określona przez
$\int_{[0,1]}|f(t)|^p\,\mathrm d\lambda(t)<1$ jest ograniczona.
Nie jest jednak normowalna, ponieważ nie jest lokalnie
wypukła (przykład~\ref{ex:quasinorm-convex-hull}).
Tak więc \emph{lokalna ograniczoność nie oznacza normowalności}.
\end{przyklad}

\subsection{Rodziny operatorów: ograniczoność punktowa i równociągłość}

\begin{definicja}[Ograniczoność punktowa i równociągłość]
\label{def:operator-family}
Niech $X,Y$ będą TVS i niech
$\mathcal T\subseteq\mathcal L(X,Y)$.
Rodzinę $\mathcal T$ nazywamy \emph{punktowo ograniczoną}, jeśli
dla każdego $x\in X$ zbiór $\{Tx:T\in\mathcal T\}$
jest ograniczony w $Y$. Rodzinę nazywamy
\emph{równociągłą} (w zerze), jeśli dla każdego otoczenia
$V$ zera w $Y$ istnieje otoczenie $U$ zera w $X$ takie,
że $T(U)\subseteq V$ dla wszystkich $T\in\mathcal T$.
\end{definicja}

\begin{twierdzenie}[Równociągłość i zbiory ograniczone]
\label{thm:equi-bounded}
Niech $X,Y$ będą TVS i $\mathcal T\subseteq\mathcal L(X,Y)$.
Jeśli $\mathcal T$ jest równociągła, to dla każdego
ograniczonego $B\subseteq X$ zbiór $\mathcal T(B)$
jest ograniczony w $Y$. W szczególności każda rodzina
równociągła jest punktowo ograniczona.
\end{twierdzenie}
\begin{proof}
Ustalmy ograniczony $B\subseteq X$ i otoczenie $V$ zera
w $Y$. Z równociągłości istnieje otoczenie $U$ zera
w $X$ z $T(U)\subseteq V$ dla wszystkich $T\in\mathcal T$.
Ponieważ $B\subseteq tU$ dla pewnego $t>0$, mamy
$T(B)\subseteq tT(U)\subseteq tV$ dla każdego $T$.
Zatem $\mathcal T(B)\subseteq tV$.
\end{proof}

\begin{przyklad}[Punktowa ograniczoność bez równociągłości]
\label{ex:c00-uniform-boundedness-failure}
Niech $X=c_{00}$ będzie przestrzenią ciągów o skończonym
nośniku z normą $\|x\|_\infty=\sup_{k\ge1}|x_k|$.
Dla $n\ge1$ określmy ciągły funkcjonał
$T_n:X\to\K$, $T_nx=nx_n$. Każdy $x\in c_{00}$
ma skończony nośnik, zatem $T_nx=0$ dla wszystkich
dostatecznie dużych $n$. Rodzina $\{T_n:n\ge1\}$
jest więc punktowo ograniczona (a nawet $T_nx\to0$
dla każdego $x$). Nie jest równociągła: zbiór
$B=\{e_n:n\ge1\}$ (gdzie $e_n$ jest $n$-tym wektorem
standardowej bazy współrzędnościowej) jest ograniczony w $X$, lecz
$T_n(e_n)=n$ i $\mathcal T(B)$ jest nieograniczony w $\K$.

Ten przykład pokazuje, że punktowa ograniczoność rodzin
operatorów \emph{nie wystarcza} do równociągłości w dowolnej
przestrzeni normowanej. Zupełność dziedziny będzie kluczowa
w późniejszym twierdzeniu Banacha--Steinhausa.
\end{przyklad}

\begin{stwierdzenie}[Rodziny na przestrzeniach unormowanych]
\label{prop:equi-normed}
Niech $X,Y$ będą przestrzeniami unormowanymi, a
$\mathcal T\subseteq\mathcal L(X,Y)$. Wtedy $\mathcal T$
jest równociągła dokładnie wtedy, gdy
$\sup_{T\in\mathcal T}\|T\|<\infty$,
gdzie $\|T\|=\sup_{\|x\|\le1}\|Tx\|$.
\end{stwierdzenie}
\begin{proof}
Jeśli rodzina jest równociągła, to dla otoczenia
$V=\{y:\|y\|<1\}$ istnieje $r>0$ takie, że
$\|x\|<r$ pociąga $\|Tx\|<1$ dla wszystkich $T$.
Dla $\|x\|\le1$ zastosujmy to do $(r/2)x$:
$\|Tx\|<2/r$. Stąd supremum norm operatorowych jest
skończone.
Odwrotnie, jeśli $M=\sup_T\|T\|<\infty$, to dla
$\varepsilon>0$ wystarczy wybrać
$U=\{x:\|x\|<\varepsilon/(M+1)\}$.
\end{proof}

\subsection{Topologie zbieżności operatorów}

\begin{definicja}[Zbieżność punktowa i jednostajna na ograniczonych]
\label{def:operator-topologies}
Niech $X,Y$ będą TVS Hausdorffa. Na przestrzeni
$\mathcal L(X,Y)$ rozważamy dwie topologie:
\begin{enumerate}[label=\textup{(\roman*)}]
\item topologię $\tau_s$ \emph{zbieżności punktowej},
wyznaczoną odwzorowaniami $T\mapsto Tx$, $x\in X$;
\item topologię $\tau_b$ \emph{zbieżności jednostajnej na
zbiorach ograniczonych}, której bazę otoczeń zera
stanowią skończone przecięcia zbiorów
\[
 \mathcal N(B,V)=\{T\in\mathcal L(X,Y):T(B)\subseteq V\},
\]
gdzie $B\subseteq X$ jest ograniczony, a $V$ jest
otoczeniem zera w $Y$.
\end{enumerate}
\end{definicja}

\begin{twierdzenie}[Struktura liniowo-topologiczna przestrzeni operatorów]
\label{thm:op-topologies}
Niech $X,Y$ będą TVS Hausdorffa. Topologie $\tau_s$ i
$\tau_b$ czynią $\mathcal L(X,Y)$ przestrzenią
liniowo-topologiczną Hausdorffa, a $\tau_b$ jest silniejsza
od $\tau_s$. Dla sieci $(T_d)$ i $T\in\mathcal L(X,Y)$
zachodzi
\begin{align*}
 T_d\xrightarrow{\tau_s}T
 &\iff T_dx\to Tx\ \text{dla każdego }x\in X,\\
 T_d\xrightarrow{\tau_b}T
 &\iff \text{dla każdego ograniczonego }B\subseteq X
 \text{ i }V\in\mathcal N_Y(0)\\
 &\hspace{18mm}\text{ostatecznie }(T_d-T)(B)\subseteq V.
\end{align*}
\end{twierdzenie}
\begin{proof}
Topologia $\tau_s$ jest topologią podprzestrzeni
otrzymaną przez zanurzenie $T\mapsto(Tx)_{x\in X}$
w produkt $Y^X$, więc jest Hausdorffa i zgodna z
działaniami liniowymi.

Dla $\tau_b$ zbiory $\mathcal N(B,V)$ zawierają zero.
Wystarczy rozważać otwarte, zrównoważone otoczenia $V$,
ponieważ tworzą one bazę otoczeń zera w $Y$.
Dla takiego $V$ można dobrać otwarte, zrównoważone
$W$ z $W+W\subseteq V$;
wtedy
$\mathcal N(B,W)+\mathcal N(B,W)\subseteq\mathcal N(B,V)$.
Przecięcia wyraża się przez skończone układy warunków.
Zbiory te są zrównoważone, jeżeli $V$ jest zrównoważony.
Dla każdego $S\in\mathcal N(B,W)$ zachodzi
$S+\mathcal N(B,W)\subseteq\mathcal N(B,V)$.
Warunki te zapewniają istnienie topologii grupowej
z podaną bazą \emph{otoczeń} zera. Nie twierdzimy, że każdy
zbiór $\mathcal N(B,V)$ jest otwarty; może nim nie być.

Sprawdźmy łączną ciągłość mnożenia przez skalar.
Ustalmy $\lambda_0\in\K$, $T_0\in\mathcal L(X,Y)$,
ograniczony $B\subseteq X$ i otoczenie $V$ zera w $Y$.
Wybierzmy zrównoważone otoczenie $W$ z $W+W\subseteq V$.
Ponieważ $T_0(B)$ jest ograniczony, dla $\lambda$
dostatecznie bliskich $\lambda_0$ mamy
$(\lambda-\lambda_0)T_0(B)\subseteq W$.
W otoczeniu $\lambda_0$ liczby $|\lambda|$ są ograniczone
przez pewne $M>0$. Dobierzmy zrównoważone otoczenie
$U$ takie, że $\lambda U\subseteq W$ dla $|\lambda|\le M$.
Jeżeli $T-T_0\in\mathcal N(B,U)$, to
\[
 (\lambda T-\lambda_0T_0)(B)
 \subseteq\lambda(T-T_0)(B)+(\lambda-\lambda_0)T_0(B)
 \subseteq W+W\subseteq V.
\]
To dowodzi ciągłości łącznej. Ponieważ każdy singleton
$\{x\}\subseteq X$ jest ograniczony, otoczenia
$\mathcal N(\{x\},V)$ należą do bazy $\tau_b$,
a zatem $\tau_b\supseteq\tau_s$. Warunek Hausdorffa
wynika z faktu, że jeżeli $T\ne0$, to dla pewnego $x$
mamy $Tx\ne0$ i można oddzielić $Tx$ od zera w $Y$.
Opisy zbieżności wynikają wprost z definicji.
\end{proof}

\begin{stwierdzenie}[Ograniczoność rodzin w topologiach operatorów]
\label{prop:op-family-bounded}
Niech $X,Y$ będą TVS Hausdorffa i
$\mathcal T\subseteq\mathcal L(X,Y)$. Wówczas:
\begin{enumerate}[label=\textup{(\roman*)}]
\item $\mathcal T$ jest ograniczona w $(\mathcal L(X,Y),\tau_s)$
dokładnie wtedy, gdy jest punktowo ograniczona;
\item $\mathcal T$ jest ograniczona w $(\mathcal L(X,Y),\tau_b)$
dokładnie wtedy, gdy dla każdego ograniczonego $B\subseteq X$
zbiór $\mathcal T(B)$ jest ograniczony w $Y$.
\end{enumerate}
\end{stwierdzenie}
\begin{proof}
W (i) sprawdzamy ograniczoność we wszystkich ewaluacjach
$T\mapsto Tx$ i korzystamy z kryterium ograniczoności
w produkcie. W (ii) z definicji ograniczoności w
$\tau_b$ dla danego $B,V$ istnieje $t>0$ takie,
że $\mathcal T\subseteq t\mathcal N(B,V)$; jest to
równoważne zawieraniu $\mathcal T(B)\subseteq tV$.
Skończone przecięcia warunków nie zmieniają tego kryterium,
ponieważ bierzemy największą z odpowiednich stałych
skalowania po uprzednim zastąpieniu otoczeń zrównoważonymi.
\end{proof}

\begin{przyklad}[Dwie różne zbieżności operatorów]
Niech $X=c_{00}$ z normą $\ell^\infty$, a $Y=\K$.
Funkcjonały $T_nx=nx_n$ spełniają
$T_n\xrightarrow{\tau_s}0$, ale nie
$T_n\xrightarrow{\tau_b}0$, ponieważ
$\{e_n:n\ge1\}$ jest ograniczony i
$\sup_{x\in\{e_n:n\ge1\}}|T_nx|=n$.
\end{przyklad}

\begin{twierdzenie}[Przypadek przestrzeni unormowanych]
\label{thm:op-b-norm}
Jeżeli $X,Y$ są przestrzeniami unormowanymi, to topologia
$\tau_b$ na $\mathcal L(X,Y)$ jest topologią normy
operatorowej. Jeżeli $X$ ma dodatkowo wymiar skończony,
również $\tau_s=\tau_b$.
\end{twierdzenie}
\begin{proof}
Kula jednostkowa $B_X=\{x:\|x\|\le1\}$ jest ograniczona.
Warunek $(T_d-T)(B_X)\subseteq\{y:\|y\|<\varepsilon\}$
pociąga $\|T_d-T\|\le\varepsilon$. Wybierając $\varepsilon/2$ zamiast $\varepsilon$ widzimy, że $\tau_b$ jest co najmniej tak silna jak topologia normy. Odwrotnie, jeśli
$B\subseteq X$ jest ograniczony, to
$B\subseteq RB_X$ dla pewnego $R>0$, więc
$\|S\|<\varepsilon/R$ implikuje $S(B)\subseteq
\{y:\|y\|<\varepsilon\}$. W ten sposób otrzymujemy
tę samą topologię.

Jeśli $\dim X=0$, to $\mathcal L(X,Y)=\{0\}$
i obie topologie są jedyną topologią tej przestrzeni.
Załóżmy więc $\dim X=m\geq1$ i wybierzmy bazę
$e_1,\ldots,e_m$. Dla każdego $x=\sum_i a_ie_i$
zachodzą oszacowania
$|a_i|\le C\|x\|$ dla pewnej wspólnej stałej $C$:
aby to uzasadnić, rozważmy zbiór
$S=\{(a_1,\ldots,a_m):\max_i|a_i|=1\}$.
Jest on zwarty w $\K^m$, a ciągła funkcja
$(a_i)\mapsto\|\sum_i a_ie_i\|$ ma na nim dodatnie
minimum $c$. Skalowanie daje
$\|\sum_i a_ie_i\|\ge c\max_i|a_i|$, co
wystarcza przy $C=1/c$.
Jeżeli $\|S e_i\|<\delta$ dla każdego $i$, to
$\|Sx\|\le mC\delta\|x\|$, a zatem ograniczenie
operatora kontroluje skończenie wiele ewaluacji.
Wynika stąd $\tau_s=\tau_b$.
\end{proof}

\subsection{Twierdzenie o metryzowalności przestrzeni liniowo-topologicznych}

\begin{definicja}[Pierwszy aksjomat przeliczalności]
Przestrzeń topologiczną $X$ nazywamy
\emph{pierwszoprzeliczalną}, jeżeli każdy punkt ma
przeliczalną bazę otoczeń. W TVS wystarcza, aby
przeliczalna baza otoczeń istniała w zerze.
\end{definicja}

\begin{twierdzenie}[Kryterium metryzowalności TVS]
\label{thm:tvs-metrizable}
Niech $X$ będzie przestrzenią liniowo-topologiczną
Hausdorffa nad $\K$. Następujące warunki są równoważne:
\begin{enumerate}[label=\textup{(\roman*)}]
\item $X$ jest metryzowalna;
\item $X$ ma przeliczalną bazę otoczeń zera;
\item topologia $X$ jest wyznaczona przez pewną
      metrykę niezmienniczą względem przesunięć,
      tj. $d(x+a,y+a)=d(x,y)$ dla wszystkich $x,y,a\in X$.
\end{enumerate}
Co więcej, metryka z (iii) wyznacza kanoniczną strukturę
jednostajną $X$, a więc zupełność tej metryki jest
równoważna zupełności TVS.
\end{twierdzenie}
\begin{proof}
Implikacje $(\mathrm{iii})\Rightarrow(\mathrm{i})\Rightarrow(\mathrm{ii})$ są oczywiste:
w przestrzeni metrycznej kule o promieniach $1/n$
tworzą przeliczalną bazę otoczeń zera.

Załóżmy (ii). Wybierzmy malejącą bazę otwartych otoczeń
zera $(W_n)_{n\ge1}$ i skonstruujmy rekurencyjnie
otwarte, symetryczne i zrównoważone otoczenia $U_n$,
$n\ge0$, takie, że
$U_0=X$, $U_n\subseteq W_n$ dla $n\ge1$ oraz
$U_{n+1}+U_{n+1}+U_{n+1}\subseteq U_n$.
Jest to możliwe dzięki ciągłości dodawania i istnieniu
zrównoważonych otoczeń zera. Mamy $\bigcap_nU_n=\{0\}$:
przeciwne założenie przeczyłoby własności Hausdorffa.

Dla $x,y\in X$ połóżmy
\[
 d(x,y)=\inf\Big\{\sum_{j=1}^{m}2^{-n_j}:\
 m\ge1,\ x=x_0,\ldots,x_m=y,\
 x_{j-1}-x_j\in U_{n_j},\ n_j\ge0\Big\}.
\]
Infimum dotyczy wszystkich skończonych łańcuchów.
Ponieważ $U_0=X$, zbiór, z którego bierzemy infimum,
jest niepusty. Symetria zbiorów $U_n$ daje
$d(x,y)=d(y,x)$, a konkatenacja łańcuchów nierówność
trójkąta. Wzór zależy tylko od różnic $x-y$, więc
$d(x+a,y+a)=d(x,y)$.

Sprawdzimy, że $d$ jest dodatnio określona i generuje
topologię $X$. Potrzebujemy następującego rachunku:
\emph{jeżeli łańcuch ma łączną wagę mniejszą niż $2^{-n}$,
$n\ge1$, to różnica jego końców należy do $U_{n-1}$.}
Dowód prowadzimy indukcyjnie względem liczby jego
odcinków. Dla jednego odcinka warunek
$2^{-n_1}<2^{-n}$ daje $n_1>n$, czyli różnica należy do
$U_{n+1}\subseteq U_{n-1}$. Dla łańcucha dłuższego
wybierzmy odcinek przecinający połowę jego łącznej
wagi. Suma wag odcinków poprzedzających ten odcinek
oraz suma wag odcinków następujących po nim są
obie mniejsze niż $2^{-(n+1)}$. Z założenia
indukcyjnego ich całkowite różnice należą do $U_n$.
Waga wyróżnionego odcinka jest potęgą $2^{-n_j}$
i jest mniejsza niż $2^{-n}$, skąd $n_j\geq n+1$;
jego różnica należy więc do $U_{n_j}\subseteq U_n$.
Różnica końców całego łańcucha leży zatem w
$U_n+U_n+U_n\subseteq U_{n-1}$.

Dla każdego $n\ge1$ wynika stąd
\[
 U_{n+1}\subseteq\{x:d(x,0)<2^{-n}\}\subseteq U_{n-1}.
\]
Rzeczywiście pierwsza inkluzja odpowiada łańcuchowi
jednoodcinkowemu, a druga -- powyższemu rachunkowi.
Jeżeli $d(x,0)=0$, to $x$ należy do każdego $U_n$,
więc $x=0$. Metryka jest zatem dodatnio określona.
Dwie inkluzje dowodzą zarazem, że jej kule stanowią
tę samą bazę otoczeń zera co $U_n$; niezmienniczość
względem przesunięć daje zgodność topologii w każdym punkcie.

Otoczenia przekątnej dla metryki $d$ są postaci
$\{(x,y):d(x,y)<r\}$. Powyższe inkluzje pokazują,
że tworzą one dokładnie tę samą strukturę jednostajną,
co otoczenia $E_U=\{(x,y):x-y\in U\}$.
Sieci Cauchy'ego obu struktur są więc identyczne.
To dowodzi ostatniej części tezy.
\end{proof}

\begin{przyklad}[Metryzowalna, zupełna, ale nienormowalna]
\label{ex:metrizable-not-normable}
Na $X=\K^{\N}$ z topologią produktową określmy
\[
 d(x,y)=\sum_{n=1}^{\infty}2^{-n}\min\{1,|x_n-y_n|\}.
\]
Szereg jest zbieżny, $d$ jest metryką niezmienniczą
względem przesunięć i opisuje zbieżność współrzędną.
Przestrzeń $X$ jest zupełna jako produkt zupełnych
przestrzeni $\K$ (co wynika z zupełności produktów),
lecz nie jest normowalna, ponieważ nie ma ograniczonego
otoczenia zera. To pokazuje, że metryzowalność i
zupełność nie zastępują normowalności.
\end{przyklad}

\subsection{Jakie założenia będą potrzebne w następnych rozdziałach?}

Różnicę między ograniczonością punktową i równociągłością
usuwa w ważnych sytuacjach zasada jednostajnej ograniczoności;
jednym z jej naturalnych kontekstów są przestrzenie beczkowe.
Twierdzenie to wymaga dodatkowej struktury i zostanie
udowodnione później, a nie założone milcząco.

W kolejnym rozdziale wprowadzimy lokalną wypukłość,
seminormy i funkcjonał Minkowskiego. Dowiemy się między
innymi, dlaczego wypukłość otoczeń pozwala zamieniać
ograniczoność na oszacowania przez seminormy oraz
kiedy lokalna ograniczoność prowadzi już do normowalności.
